% Research manuscript. AC-G139 identifiers reserved; not installed.
% The G138 dependency is kept at its exact completed-review scope.
\section{Conjugate quartic characters and two quadratic compensations}
\label{sec:cubic-quartic-compensated-families}

Put \(p=269\), \(K=\mathbb Q(\sqrt{-3})\), and
\(\eta=\chi_{-3}\).  Let \(\psi_5\) be the primitive quartic
character with \(\psi_5(2)=i\), where \(i\equiv82\pmod p\),
and let \(\Psi\) be its arithmetic Galois avatar.
Let \(\tau_8\) and \(\tau_p\) be the quadratic characters of
\(\mathbb Q(\sqrt2)\) and \(\mathbb Q(\sqrt p)\), respectively.
The same notation will denote their Galois avatars.  In particular
\(\tau_p\) is the Teichmuller lift of \(\bar c_p^{134}\),
where \(c_p\) is the actual cyclotomic action.
For \(\sigma\in\{1,-1\}\) and \(a,b\in\{0,1\}\), define
\[
 \psi_{\sigma,a,b}=\psi_5^\sigma\tau_8^a\tau_p^b,
 \qquad N_{a,b}=10\cdot4^a\cdot269^b.
 \tag{QCF.1}
\]
The character in (QCF.1) is regarded as a character of the displayed
even modulus.  Its primitive conductor is \(5\cdot8^a\cdot269^b\).
For every occurrence of an original model retain its own original
\(\mathcal O_0\), \(\mathrm{GD}/\mathrm{Dat}\), marked data,
selected local prime \(P_{*,\sigma,a,b,e}\), and complete model.
The old modularity premise still includes \(e_{\mathcal O_0}=1\).

The weights considered here are
\[
 e>0,\qquad e\equiv133\ \hbox{or}\ 267\pmod{268}.
 \tag{QCF.2}
\]
Define \(r\in\{133,267\}\) by
\[
                         r\equiv e+134b\pmod{268}.
 \tag{QCF.3}
\]
Thus \(b=1\) exchanges the two residues.  The auxiliary host
construction used for \(r=267\) is AC-G136, and for \(r=133\)
it is AC-G138.  All their named native, published, class-field,
cohomological, and generic algebraic inputs are retained, with the
class inputs exactly AC-G135 and AC-G137, respectively.  In using
AC-G138, its complete proof and reviewed scope are a dependency;
an unfinished draft is not an additional assumed theorem.

\begin{lemma}[AC-G139-TypeII: exact quadratic reduction of every endpoint]
\label{AC-G139-TypeII}
Retain the original modularity and local hypotheses of AC-G68-point
and the exact published attachment, local compatibility, primitive
newform, and Katz inputs of AC-G128.  Then
there is no characteristic-zero Type II point for any original
model in (QCF.1)--(QCF.2), for either cubic label \(j=1,2\).
This assertion does not require a host or an auxiliary seed.
\end{lemma}
\begin{proof}
The character \(\psi_{\sigma,a,b}\) is odd of order four.
Both quadratic factors are even, and
\(\operatorname{comp}_p\psi_{\sigma,a,b}=\tau_p^b\).
The complete literal \(\mathrm{GD}\) conditions hold in the same
specified source domain: the parity equality is
\(\psi(-1)=(-1)^e=-1\), \(268\nmid2e\),
\(3\nmid N_{a,b}\), and \(p\nmid3h_K\).
The \(\mathrm{GD.eta}\) kernel condition follows from inertia
at three: the cyclotomic and prime-to-\(p\) character fields are
unramified there, whereas \(K/\mathbb Q\) is ramified.
There is no \(p\nmid N\) clause in the literal definition.

AC-G68 supplies an actual odd irreducible integral modular
representation \(V\) of \(G_{\mathbb Q}\), with
\[
 \begin{gathered}
 \det V=\eta,\quad R_f=Vc_p^{-e},\quad
 \operatorname{wt}(f)=2e+1,\\
 \bar V^{\rm ss}=
  \eta\bar\Psi^\sigma\bar\tau_8^a\bar\tau_p^b\bar c_p^e
   \ \oplus\ 
  \bar\Psi^{-\sigma}\bar\tau_8^a\bar\tau_p^b\bar c_p^{-e}.
 \end{gathered}
 \tag{QCF.4}
\]
All original coefficient and selected-point hypotheses are used
before the following operation.  Put
\[
 \zeta=\Psi^{1-\sigma}\tau_8^a\tau_p^b,
                      \qquad W=V\zeta.
 \tag{QCF.5}
\]
Here \(\Psi^{1-\sigma}\) is one or the quadratic character
\(\chi_5=\Psi^2\).  Thus \(\zeta^2=1\), its values are in
\(\{1,-1\}\), and neither the original lattice nor coefficient
ramification index is changed.  Since
\(1-2\sigma\equiv-1\pmod4\), (QCF.4) gives
\[
       \det W=\eta,\qquad
       \bar W^{\rm ss}=
        \eta\bar\Psi\bar c_p^e\oplus
                         \bar\Psi^{-1}\bar c_p^{-e}.
 \tag{QCF.6}
\]
The actual primitive automorphic quadratic twist of \(f\)
has attached representation \(Wc_p^{-e}\).  Quadraticity
removes the inverse-character ambiguity in the geometric-Frobenius
attachment convention.  This is a twist of the actual cuspidal
representation, not an imprimitive coefficientwise twist.

Away from \(p\), the exact residual-kernel condition gives
finite inertia \(\psi_{\sigma,a,b}\oplus
\psi_{\sigma,a,b}^{-1}\) on \(G_K\), not merely this
semisimplified reduction.  Tensoring by (QCF.5) gives
\(\Psi\oplus\Psi^{-1}\).  At rational three, inertia
over \(K_3\) is trivial and determinant \(\eta\) forces
the characters \(1,\eta\).  The quadratic twist is unramified
at three.  Consequently \(W\) is unramified at every rational
prime outside \(3,5,p\); its inertia at five is
\(\Psi\oplus\Psi^{-1}\), with zero monodromy everywhere away
from \(p\).  In particular the ramified \(\tau_8\) at two
has actually been removed.

At \(p\), use this original model's own AC-G65 flag and the
exact AC-G68 untwisting.  The resulting \(V|_{G_{K_p}}\)
has potentially crystalline type \(\tau_p^b\oplus\tau_p^b\)
and ordered inertia \(\tau_p^bc_p^e,\tau_p^bc_p^{-e}\).
Both \(\Psi\) and \(\tau_8\) are unramified at \(p\).
Thus \(W|_{G_{K_p}}\) has trivial potentially crystalline
type.  Choose a finite Galois extension on which it is crystalline.
Inertia acts trivially on its full crystalline period module;
unramified Galois descent of that module proves crystallinity
already over \(K_p\).  As \(K_p/\mathbb Q_p\) is unramified,
the same semilinear descent proves crystallinity over
\(\mathbb Q_p\).  The ordered line descends to the rational
decomposition group: its inertia character \(c_p^e\) is
different from \(c_p^{-e}\) on every open inertia subgroup,
so its labelled line is unique and is preserved by conjugation.
Its quotient is \(\eta\delta^{-1}\) if the first character
is \(\delta\); the unramified factor \(\eta\) is retained.
This is the good-twist local argument with an unramified determinant,
not a literal application of the split-prime statement
AC-K29-good-twist at the inert prime \(269\).

All hypotheses used in the representation-level clause of
AC-G128-quartic-exclusion now hold for \(W\).  Its primitive
modular level is \(75\), and the primitive \(\eta\psi_5\)
twist has level \(15\), as in that proof.  The two weights in
(QCF.2) give the same weight residue \(267\) modulo \(268\)
for this level-fifteen form.  The already checked nonzero constant
\(145\) and Katz contradiction therefore exclude \(W\), hence
\(V\).  No map from the original model in (QCF.1) to the
original \(N=10\) selected model has been asserted or needed.
\end{proof}

\begin{theorem}[AC-G139-families: fresh support and the original compensated models]
\label{AC-G139-families}
Retain the exact completed host constructions and inputs specified
above.  For each \(\sigma,r,j\), there are a finite CM host,
actual normalized tuple rings \(B\subset A\subset R\),
new native maps \(\Phi_2,\Phi_3\), and proper native modular
ideals \(J_2,J_3\), for the residual tuple
\[
 \theta_{\sigma,r}=
  (\bar\chi^j\bar\Psi^\sigma\bar c_p^r,
                  1,\bar\chi^j\bar\Psi^{-\sigma}\bar c_p^{-r}).
 \tag{QCF.7}
\]
They supply a uniform even ideal \(M_E\) and, with
\(I=\Phi_3(J_3)B\), \(I_{\rm off}=I\cap M_E\), satisfy
\[
 \begin{aligned}
 I&\subset\mathfrak q\cap B
                 &&(\mathfrak q\text{ Type I}),\\
 I_{\rm off}&\subset\mathfrak q\cap B
                 &&(\mathfrak q\text{ off triangular}),\\
 \sqrt{I_{\rm off}+\lambda_\theta}&=\mathfrak m_B.
 \end{aligned}
 \tag{QCF.8}
\]
For fixed \(\sigma,r,j,\mathcal O_0\), these host data may
be chosen once for all \(a,b,e\) satisfying (QCF.2)--(QCF.3).
Every such original model satisfies all five literal
\(\mathrm{M1}\)--\(\mathrm{M5}\), has finite complete degree-zero
ring over its original \(\mathcal O_0\), and has only rigid
Type I two-sided characteristic-zero points.

For each original row separately, let \(L_j\) be its original
weight-minus-one cotangent length, and let \(s_j\) be its
normalized strict Selmer length for
\[
       (E_0/\mathcal O_0)
        (\chi^{-j}\Psi^{-\sigma}\tau_8^a\tau_p^bc_p^{-e}).
 \tag{QCF.9}
\]
Put \(Y=N_{a,b}(e+1)\) and \(b_e=121^e+1+121^{-e}\).
For every nonempty \(J_0\subset\{1,2\}\), one nonzero rational
integer \(\mathcal N\) gives
\[
 \begin{gathered}
 s_j\le L_j/e_{\mathcal O_0}<\infty,\qquad
 \sum_{j\in J_0}L_j/e_{\mathcal O_0}\le2v_{269}(\mathcal N),\\
 \log|\mathcal N|\le5C_1\#J_0Y^{A_1+1}\log11,\qquad
 \left(\sum_{j\in J_0}L_j/e_{\mathcal O_0}\right)\log269
       \le10C_1\#J_0Y^{A_1+1}\log11.
\end{gathered}
 \tag{QCF.10}
\]
The constants and original-field classifier are those of AC-G112.
These are finitely specified character families; no all-character,
all-coefficient-prime, or H04.4 family-budget conclusion is claimed.
\end{theorem}

\begin{proof}
\proofstep{1. A fresh conjugate auxiliary branch.}
The auxiliary prime is \(q_{\rm aux}=16141\), with its fixed
order-\(p\) character \(\theta_0\) and real cyclic field \(T\).
For \(r=267\) use the actual integral cusp eigenlift of AC-G136;
for \(r=133\) use the independently constructed eigenlift of
AC-G138.  In both cases it is
\[
       g\in S_{2r+1}(15q_{\rm aux},\xi_0\theta_0),
       \qquad \xi_0=\eta\psi_5^{-2}=\chi_{-15}.
 \tag{QCF.11}
\]
Their integral descent and full Hecke eigenlifting, including the
bad and ordinary roots, remain exactly those proved there.  They
do not depend on the sign \(\sigma\), since
\(\psi_5^{-2\sigma}=\psi_5^{-2}\).
Define actual primitive twists
\[
 \begin{split}
 f_{*,\sigma}&=g\otimes
                 (\eta\psi_5^{-\sigma}\theta_0^{-1/2}),\\
 V_\sigma&=R_g\eta\Psi^\sigma\Theta_0^{1/2}c_p^r.
 \end{split}
 \tag{QCF.12}
\]
The inverse automorphic/Galois character convention is the one
checked in AC-G118 and AC-G136.  Since
\(\det R_g=\eta\Psi^2\Theta_0^{-1}c_p^{-2r}\),
\[
 \det V_\sigma=\eta,\qquad
 \bar V_\sigma^{\rm ss}=
     \eta\bar\Psi^\sigma\bar c_p^r
           \oplus\bar\Psi^{-\sigma}\bar c_p^{-r}.
 \tag{QCF.13}
\]
Each primitive level is exactly \(75q_{\rm aux}^2\), with
nebentype \(\eta\).  The inertia pairs are
\(1,\eta\) at three, \(\Psi^\sigma,\Psi^{-\sigma}\) at five,
\(\Theta_0^{1/2},\Theta_0^{-1/2}\) at \(q_{\rm aux}\),
and trivial elsewhere away from \(p\).  At \(p\) it is
crystalline ordinary.  In particular the new branch is unramified
at two and eleven.  Notice that \(V_{-1}=V_1\chi_5\), but
the later rank-three representation is not obtained by a scalar
twist of the whole original rank-three seed.

Choose an actual reduced domain ordinary branch through each
\(f_{*,\sigma}\).  Its dense weights are \(w\equiv r\pmod{268}\).
The characteristic-zero tame character selection is a domain
quotient, not an integral projector dividing by \(p\).
The same actual specialization/control inputs fix its full old
inertia.  Impose the corresponding old compact type at five
with its newly labelled characters; do not identify its residual
Hecke ideal with the one for the other sign.

\proofstep{2. The exact-root and finite-type construction for this branch.}
Start from \(F_1=K\mathbb Q(\sqrt{-11})\) and \(F_*=F_1T\).
The field and local checks of AC-G136 and AC-G138 are unchanged:
the CM extension is finite-unramified, relatively split at every
ramified old place and at \(p\), and its \(p\)-completions are
unramified of degree \(538\).  The elementary irreducibility
checks apply to \(V_\sigma\): the quadratic restriction would
give a self-twist ramified at eleven, and the later degree-\(p\)
restriction cannot give a rank-two character self-twist of order
\(p\).  The ordinary inert-prime exclusion over \(K\) is
unchanged by the quadratic scalar twist.

Set \(E=\mathbb Q(\sqrt2)\), \(u=(1+\sqrt2)^2\), and
\(M=E(\mu_p,u^{1/p})\).  The exact elements \(h\in I_2\)
and \(t\in I_p\), and the subsequent conjugation of \(h\),
can be used in the new branch's own exact-root argument.
Both have \(\chi_5=1\): it is unramified on both inertia
groups.  More directly, this branch is unramified at two and
its inertia trace at \(t\) is exactly
\(-((1+W)+(1+W)^{-1})\), since every \(w\) is odd.
The single-element construction therefore supplies a new prime
\(q\), if necessary depending on \(\sigma,r\), with
\[
 q\equiv-1\pmod p,\quad v_p(q+1)=1,\quad
 \chi(\operatorname{Fr}_q)=\Psi(\operatorname{Fr}_q)=1,
 \quad\eta_E(q)=-1,
 \quad u\in(\mathbb F_{q^2}^{\times})^p.
 \tag{QCF.14}
\]
No common prime for different branches is required.

Use the full unit-torus finite character obtained from
\((\chi^j,1,\chi^j)w_0(r-1,0,1-r)^{-1}\).
It is the same formula for both signs, because \(\Psi\) is
unramified at \(p\).  Its constancy on the dense progression,
the normalized ordinary valuation \(2f-2fw+2f(w-1)=0\),
and strict rearrangement for the other orders remain literal.
The integral old and finite-torus projectors commute with the
inverse traces and the finite group at \(q\).  The residual
generic witness has \(\Psi=1\), so the elliptic tests and the
projective finite-type positivity proof apply to this new
localization without change of their hypotheses.
More explicitly, take the actual finite free inverse-trace ordinary
module with these projectors and its finite Hecke image.  The
dense classical transfers from the new branch satisfy every
operator relation in that image.  The AC-G78 Cayley--Hamilton
specialization argument gives its character at the exact-root
point, and one localizes at that character's own residual maximal
ideal.  This is a newly obtained nonzero Hecke localization for
\(\theta_{\sigma,r}\), not the old sign's localization.
The actual finite free multiplicity module then supplies the
type at the classical weight \(w=r\).  The ray cover of \(E\)
kills its specified tame parameter at \(q\), and gives an actual
rank-three ordinary cuspidal seed.  Its residual determinant is
the one in (QCF.7); its old type at five has the new labelled
characters.  This constructs the required support through
actual modular cohomology, rather than transporting a putative
rank-three scalar twist or a formal characteristic polynomial.
The rank-two seed also stays irreducible after adjoining \(E\):
a new quadratic self-twist would be ramified at two, where the
representation is unramified.  The subsequent cyclic
\(p\)-power restrictions preserve both seeds by the same
determinant and full ordered-flag arguments as in AC-G136.

\proofstep{3. Class sectors, filling, and both native components.}
The maximal prime-to-\(p\) host is still
\(H_{\rm ab}=K\mathbb Q(\sqrt{-11})\mathbb Q(\sqrt2)\).
For \(r=267\), its four strict adjacent sectors use
\(d=1,-11,2,-22\); for \(r=133\), they use
\(d=p,-11p,2p,-22p\).  Replacing \(\Psi\) by
\(\Psi^{-1}\) does not change any of these splitting fields:
\(\ker(\Psi\tau_D)=\ker(\Psi^{-1}\tau_D)\), because the
two characters are inverses.  The target Kummer eigenspace is
exchanged with the other quartic eigenspace.  The class theorems
kill every nontrivial cubic class sector, and the equivariant
unit character gives dimension at most one in either quartic
eigenspace.  The base strict-zero inputs AC-G130 and AC-G134
explicitly cover both quartic signs and both cubic labels.
Consequently the same bound \(0+1+1+1<4\) holds.

The cyclic \(p\)-ascent, selected-place nonzero classes and
filling therefore work with the new tuple.  For their residual
disjointness step use the splitting field generated by
\(K,\chi,\Psi,\bar c_p\), which is independent of \(\sigma\)
and unramified at two and eleven.  Do not substitute the enlarged
original splitting field containing \(\tau_8\) into that step.
The dual characters controlling the mixed filling obstruction
depend on \(\Psi^2\), so are unchanged.  The other rational
characters are included with their new sign in the selected-place
splitting cover; they remain nontrivial since
\(\bar\Psi^\sigma\bar c_p^r(c)=1\) and \(\eta(c)=-1\).
For both residues, \(2r\equiv266\pmod{268}\);
the actual local endpoint square has order \(134\), and the
dual local exponents are \(267\) and \(3\).  The actual
\(\mathrm{Psi_N}\), finite class exponent, generator, Euler,
and rational-frame hypotheses are therefore the same checked
ones.  At a sufficiently high cyclotomic layer one has
\(h+a+z+1<d\) and every \(d_v>5\).  The scalar auxiliary
place can again use \(\Psi=-1\), unchanged by inversion.

Construct the rank-two native Pair on this rational tuple with
the actual \(\kappa_2\), and its arithmetic point from
\(V_\sigma\kappa_2\).  The weights
\(b_\kappa(\tau)-r,b_\kappa(\tau)+r\), finite phases,
away kernels and auxiliary unramifiedness satisfy the same
literal PT fields.  The new trace/projector map \(\Phi_2\)
is surjective, its new \(J_2\) is proper, and the complete
rank-two component theorem gives uniform even support, including
the independent middle-character factor.

For rank three use \(\kappa_3=c_p^{-1}\), not \(\kappa_2\).
The seed from Step 2 supplies PT and makes its newly constructed
\(J_3\) proper.  The determinant/projector surjection
\(\Phi_3\), component recovery \(\beta\Phi_3=q\), and
joint-kernel FinLam argument all apply to (QCF.7).
The normalized collision, local presentation, parity-chain and
filled-path bounds now supply every input to AC-G108.  This proves
(QCF.8) for the new localization.  In particular the ideal used
for Type I points is \(I\), not merely \(I_{\rm off}\).
This construction supplies the conjugate quartic native support;
it does not assume an unnamed automorphy or OffFin statement.

\proofstep{4. Each original row reaches this host through its own flags.}
The host contains \(E=\mathbb Q(\sqrt2)\), so \(\tau_8\)
is trivial on the entire host group and all its full local groups.
Equation (QCF.3) gives the exact residual identity
\[
 (\bar\Psi^\sigma\bar\tau_8^a\bar\tau_p^b\bar c_p^e)|_{G_H}
                  =(\bar\Psi^\sigma\bar c_p^r)|_{G_H}.
 \tag{QCF.15}
\]
Away from \(p\), \(\tau_p\) is unramified and \(\tau_8\)
is killed on the host, so the original modified Kill condition
has exactly the required joint inertia kernel.  This includes
the newly ramified rational place two when \(a=1\).

At \(p\), apply AC-G65 to the specific
\(P_{*,\sigma,a,b,e}\), not a component imported from the
base row.  Its hypotheses hold, and its all-test-algebra flag is
\[
 \chi^j\Psi^\sigma\tau_8^a\tau_p^bc_p^e,
       \quad1,\quad
 \chi^j\Psi^{-\sigma}\tau_8^a\tau_p^bc_p^{-e}.
 \tag{QCF.16}
\]
The smooth inertia type is \((\chi^j\tau_p^b,1,\chi^j\tau_p^b)\).
The unramified value of \(\tau_8\) need not be one over
\(\mathbb Q_p\); it becomes one on the actual host because
that host contains its quadratic field.  The order-two ramified
factor \(\tau_p\) is not discarded.

On the full finite unit torus, comparison with Step 2 gives
the endpoint exponents
\[
                     134b+(e-r),\qquad134b-(e-r).
 \tag{QCF.17}
\]
Both are divisible by \(268\).  The parallel algebraic action
on residue units factors through the norm to \(\mathbb F_p^\times\);
the quadratic character is the same norm to exponent \(134b\).
This verifies the actual finite projector, including the
Teichmuller part.  The principal-unit algebraic weights remain
variable, and the full characteristic-zero characters for different
rows are not being identified.

For an original characteristic-\(p\) DVR curve, choose a common
finite residue extension \(k''\), and use the actual host test
ring \(k+z k''[[z]]\).  The paired labelled Hensel projectors
and strict normalization descend (QCF.16) to that ring.  The
actual universal chart gives a map annihilating \(\lambda_\theta\).
Equation (QCF.8), the balanced-monomial argument and the torsion-free
principal units of this characteristic-\(p\) domain force the
host characteristic polynomials to be their seed values.
Every element of the original \(G_K\) has a positive power in
\(G_H\), so its characteristic coefficients are algebraic
constants fixed by reduction.  This works even when \(p\)
divides the host index.  The exact original integral finite
criterion gives M2, including vertical components.

For an original characteristic-zero Type I point, use the
inverse image of the host residue field in a common finite
value ring.  Its own (QCF.16), kernels and normalization give
an actual point of \(A\).  The all-open irreducibility argument
applies because the ratios on open inertia have nonzero exponents
\(\pm e,\pm2e\).  The native point kills \(J_3\) by the first
line of (QCF.8).  Its finite phase is retained in
\(\mathrm{IsWtPart.ph}\), with weights \((1-e,1,1+e)\).
Here \(e\ge133\), so no distinct ratio is \(c_p\).
The native end2 direction gives actual automorphy and host
adjoint rigidity, and restriction/corestriction gives the
original rational adjoint rigidity.  No comparison between
the original rings in different rows has been used.

\proofstep{5. Original model conditions and the original-level count.}
The lemma excludes every original Type II endpoint in these
rows.  Thus every two-sided original point is Type I and rigid.
The exact normalized two-derivation argument gives M4;
AC-G66 gives M1 and M3, and Step 4 gave M2.
At the degree-two place over eleven, every quadratic factor
has value one, \(\Psi(11)=1\), and \(\chi\) has value one.
This prime is good for every \(3pN_{a,b}\).  The original
seed trace is \(b_e\), while every conjugate of an actual
Type I trace has absolute value at most three.  The original
degree-zero trace difference is therefore the required M5 function.

The actual host tower still consists only of quadratic and
cyclic \(p\)-power steps, so the original-field descent and
WitA/unlinking inputs apply.  AC-G112-local uses the original
finite character of conductor dividing \(N_{a,b}\), allows
\(p\mid N_{a,b}\), and gives principal compact level
\(269N_{a,b}\).  No auxiliary prime or host conductor enters
that level.  Apply its genuine \(\mathrm{Af.TrData}\) carrier
with bound three and cardinality \(C_1Y^{A_1}\), followed by
the original LemmaM and strict LemmaM' injection.  This gives
(QCF.10) exactly as in AC-G136, with
\(v_{\mathcal O_0}=e_{\mathcal O_0}v_{269}\).
Host coefficient ramification does not rescale the original lengths.
\end{proof}

\paragraph{Scope and precise boundaries.}
There are eight characters in (QCF.1) and two allowed weight
residues; the effective-residue rule specifies the host construction
for each of these sixteen combinations.  Separate hosts may be used
for different \(\sigma,r,j\).  The original moduli are
\(10,40,2690,10760\), with principal counting levels
\(2690,10760,723610,2894440\), respectively.
If an actual endpoint existed before the contradiction, its primitive
GL2 level would be \(75\cdot64^a\cdot269^{2b}\): the two
characters at two each have conductor exponent three, and the two
potentially crystalline Weil characters at \(p\) each have tame
conductor exponent one when \(b=1\).  This primitive level is
distinct from the original formal source parameter \(3N_{a,b}\).

The base strict-zero theorems do not by themselves prove original
strict zero in the \(\tau_8\)-compensated rows.  This proof supplies
their model and finite counting conclusions through actual host
support instead.  It supplies no effective first cyclotomic layer,
all-character average, or positive spectral gap.
The previously audited finite-cover/local-system dictionary applies
to the actual restriction and character twists; no new motive or
cohomological comparison is being asserted.  No Lean run is claimed.
